\documentclass{article}

\usepackage{amsmath,amssymb}
\usepackage{xurl}
\usepackage[hidelinks]{hyperref}
\setlength{\emergencystretch}{2em}

% Shared commands for notes in docs/paper-gaps/.

\DeclareUnicodeCharacter{2080}{\ifmmode{}_{0}\else\textsubscript{0}\fi}
\DeclareUnicodeCharacter{2081}{\ifmmode{}_{1}\else\textsubscript{1}\fi}
\DeclareUnicodeCharacter{2082}{\ifmmode{}_{2}\else\textsubscript{2}\fi}
\DeclareUnicodeCharacter{2099}{\ifmmode{}_{n}\else\textsubscript{n}\fi}

\newcommand{\Lean}{\textsc{Lean}}
\ExplSyntaxOn
\NewDocumentCommand{\leanid}{m}{
  \ifmmode
    \textsf{\detokenize{#1}}
  \else
    \group_begin:
    \urlstyle{sf}
    \tl_set:Nn \l_tmpa_tl {#1}
    \tl_replace_all:Nnn \l_tmpa_tl {\_} {_}
    \exp_args:NV \path \l_tmpa_tl
    \group_end:
  \fi
}
\ExplSyntaxOff
\providecommand{\tr}{\operatorname{tr}}
\newcommand{\tnleanIssueBase}{https://github.com/LionSR/TNLean/issues/}
\newcommand{\tnleanPrBase}{https://github.com/LionSR/TNLean/pull/}
\newcommand{\tnleanDocBase}{https://sirui-lu.com/TNLean/docs/}
\newcommand{\ghissue}[1]{\href{\tnleanIssueBase#1}{GitHub issue \##1}}
\newcommand{\ghpr}[1]{\href{\tnleanPrBase#1}{PR \##1}}
\newcommand{\doclink}[1]{\href{\tnleanDocBase#1.html}{\path{#1}}}

% Dirac notation and the identity operator, as in blueprint/src/macros/common.tex.
% \providecommand keeps note-local definitions authoritative.
\usepackage{dsfont}
\providecommand{\ket}[1]{|#1\rangle}
\providecommand{\bra}[1]{\langle #1|}
\providecommand{\braket}[2]{\langle #1 | #2 \rangle}
\providecommand{\ketbra}[2]{|#1\rangle\!\langle #2|}
\providecommand{\Id}{\mathds{1}}
\providecommand{\one}{\mathds{1}}
\providecommand{\C}{\mathbb{C}}
\providecommand{\MN}[1]{M_{#1}(\C)}
\providecommand{\MD}{M_D(\C)}

% External referencing for paper sources.  The build helper writes sanitized
% auxiliary files under build/paper-aux/ and excludes duplicated source labels.
\usepackage{xr}

% Import a generated paper auxiliary file with a caller-supplied label prefix.
% The candidate paths support builds from docs/paper-gaps/, the repository root,
% and build/paper-aux/ itself.
\newcommand{\importpaperaux}[2]{%
  \IfFileExists{../../build/paper-aux/#2.aux}{%
    \externaldocument[#1]{../../build/paper-aux/#2}%
  }{%
    \IfFileExists{build/paper-aux/#2.aux}{%
      \externaldocument[#1]{build/paper-aux/#2}%
    }{%
      \IfFileExists{#2.aux}{%
        \externaldocument[#1]{#2}%
      }{}%
    }%
  }%
}

% General external reference macros
\newcommand{\paperref}[2]{\ref{#1-#2}}
\newcommand{\papereqref}[2]{(\ref{#1-#2})}

% CPSV16 source (1606.00608 / MPDO-22-12-17-2.tex) external document and shortcuts
\importpaperaux{cpsv16-}{1606.00608/MPDO-22-12-17-2-tnlean}
\newcommand{\cpsvref}[1]{\paperref{cpsv16}{#1}}
\newcommand{\cpsveqref}[1]{\papereqref{cpsv16}{#1}}

% Verdict marker: \gapnote{<kind>}{<status>}. Kind is the policy
% classification (clarification, local-correction, scope-restriction,
% unfaithful, false-source, open-gap); status is open, wip (elimination
% underway), resolved, or historical. Machine-read by the site index;
% prints a verdict line.
\newcommand{\gapnote}[2]{\par\noindent\textbf{Verdict:} #1 (#2).\par}


\newcommand{\MPV}[2]{V^{(#1)}\!\left(#2\right)}

\title{Removing Length-Zero Bookkeeping from the\\
  Blocked Canonical-Form Decomposition}
\author{}
\date{2026-05-30}

\begin{document}
\maketitle
\gapnote{clarification}{resolved}

\noindent
This note explains why an all-length matrix-product-vector identity is not the
appropriate formal statement of the blocked canonical-form decomposition. It
also records the removal of that formulation from the physical canonical-form
statements.\footnote{%
  Context: PRs \ghpr{2153}, \ghpr{2156}, and \ghpr{2167} removed roughly a
  thousand lines of ``zero-tail'' plumbing built on the all-length form
  discussed here. This note records why that form is a non-source-faithful
  overcomplication and why it has now been removed.}

\section{The decomposition}

Let $A$ be a translation-invariant matrix-product tensor of bond dimension
$D$, with one-site matrices $A^i\in M_D(\mathbb C)$. For a length $N\geq0$
and a word $\sigma=(i_1,\ldots,i_N)$, define the corresponding
matrix-product-vector coefficient by
\begin{align}
    \MPV{N}{A}_\sigma
        &= \tr\!\left(A^{i_1}\cdots A^{i_N}\right).
    \label{eq:cpsv16_zero_tail_length_zero_decomposition_mpv_coefficient}
\end{align}
For $N=0$, the word is empty and the product in
Eq.~\ref{eq:cpsv16_zero_tail_length_zero_decomposition_mpv_coefficient} is the
identity on the bond space.

The blocked canonical-form reduction in Section~2.3 of
Ref.~\cite{Cirac2016MPDO_arXiv}, following Theorem~3 of
Ref.~\cite{PerezGarcia2007MPSReps}, states that after blocking by some period
$p\geq1$, the tensor is gauge equivalent to the direct sum of an all-zero
block and a weighted family of normal blocks:
\begin{align}
    A^{[p]}
        &\cong \mathbf 0_z\oplus\bigoplus_k\mu_k B_k.
    \label{eq:cpsv16_zero_tail_length_zero_decomposition_block_decomposition}
\end{align}
Here $\mathbf 0_z$ is the zero tensor on a bond space of dimension $z$; each
$B_k$ is a normal block of bond dimension $D_k$, meaning primitive,
trace-preserving, and irreducible; and each weight satisfies $\mu_k\neq0$.
The bond dimensions obey
\begin{align}
    z+\sum_k D_k
        &= D.
    \label{eq:cpsv16_zero_tail_length_zero_decomposition_bond_dimension}
\end{align}

\section{The length-zero issue}

The direct sum in
Eq.~\ref{eq:cpsv16_zero_tail_length_zero_decomposition_block_decomposition}
gives, for every length $N$ and word $\sigma$,
\begin{align}
    \MPV{N}{A^{[p]}}_\sigma
        &= \MPV{N}{\mathbf 0_z}_\sigma
           +\MPV{N}{\textstyle\bigoplus_k\mu_k B_k}_\sigma.
    \label{eq:cpsv16_zero_tail_length_zero_decomposition_all_length_identity}
\end{align}
The zero block contributes nothing at positive length. Indeed, every product
of one or more tensor matrices restricted to the zero summand vanishes, so
\begin{align}
    \MPV{N}{\mathbf 0_z}_\sigma
        &= 0,
        \qquad N\geq1.
    \label{eq:cpsv16_zero_tail_length_zero_decomposition_zero_block_positive_length}
\end{align}
At length zero, the unique coefficient is the trace of the identity on the
bond space. For any tensor $C$ with bond dimension $D_C$,
\begin{align}
    \MPV{0}{C}_{\varnothing}
        &= \tr(\mathbf 1_{D_C}),
    \label{eq:cpsv16_zero_tail_length_zero_decomposition_empty_word_trace}\\
    \MPV{0}{C}_{\varnothing}
        &= D_C.
    \label{eq:cpsv16_zero_tail_length_zero_decomposition_empty_word_dimension}
\end{align}
Consequently, the length-zero instance of
Eq.~\ref{eq:cpsv16_zero_tail_length_zero_decomposition_all_length_identity}
contains only the bond-dimension equality in
Eq.~\ref{eq:cpsv16_zero_tail_length_zero_decomposition_bond_dimension}.
Every positive length instead carries the substantive weighted normal-block
comparison.

The all-length formulation couples this elementary dimension identity to the
physical positive-length comparison. It then forces both statements to pass
together through blocking, reindexing, and sector grouping. In the formal
development, this coupling produced transport lemmas used only to move the
empty-word instance of
Eq.~\ref{eq:cpsv16_zero_tail_length_zero_decomposition_all_length_identity}
through those operations.\footnote{%
  These were the \leanid{zeroTail_*} lemmas in
  \path{TNLean/MPS/CanonicalForm/SectorComparison/}; they are removed in
  PRs \ghpr{2156} and \ghpr{2167} once the comparison is stated at positive
  length.}

The all-length formulation is also not source-faithful. The article introduces
the zero block once to account for the bond dimension and then discards it,
comparing only the nonzero normal blocks in the thermodynamic limit
\cite[Section~2.3]{Cirac2016MPDO_arXiv}. The summand $\mathbf 0_z$ preserves a
common bond space in the formal decomposition, but it represents no state and
has no normal-block structure.

\section{The positive-length statement}

The dimension statement and the physical comparison should be separated. For
every $N\geq1$, the physical comparison is
\begin{align}
    \MPV{N}{A^{[p]}}_\sigma
        &= \MPV{N}{\textstyle\bigoplus_k\mu_k B_k}_\sigma,
        \qquad N\geq1.
    \label{eq:cpsv16_zero_tail_length_zero_decomposition_positive_length_comparison}
\end{align}
Equation~\ref{eq:cpsv16_zero_tail_length_zero_decomposition_positive_length_comparison}
contains all physical content of the matrix-product-vector comparison.

Direct-sum induction independently gives the structural bond-dimension
identity for the unfiltered irreducible blocks. After discarding the all-zero
blocks, the downstream argument needs only
\begin{align}
    \sum_k D_k
        &\leq D.
    \label{eq:cpsv16_zero_tail_length_zero_decomposition_dimension_bound}
\end{align}
No physical theorem therefore needs an empty-word coefficient or a scalar
recording the dimension of the discarded bond space. Formally, positive-length
equality is \leanid{SameMPV₂Pos}; the canonical-form reduction carries
Eq.~\ref{eq:cpsv16_zero_tail_length_zero_decomposition_positive_length_comparison}
and the bound in
Eq.~\ref{eq:cpsv16_zero_tail_length_zero_decomposition_dimension_bound}.

\section{Formal status}

The declaration
\leanid{exists_tp_primitive_blockDecomp_after_blocking}, corresponding to
blueprint item \path{thm:tp_primitive_blockdecomp} in
\path{blueprint/src/chapter/ch09_canonical.tex}, now has the source-faithful
shape. Its matrix-product-vector clause is the positive-length comparison in
Eq.~\ref{eq:cpsv16_zero_tail_length_zero_decomposition_positive_length_comparison},
and its dimension clause is
Eq.~\ref{eq:cpsv16_zero_tail_length_zero_decomposition_dimension_bound}.

The intermediate declarations carrying \texttt{zeroTailDim} and the former
support lemma \leanid{zeroTail_toTensorFromBlocks_blockPower} have been
removed. The subsequent common-sector theorems now assume
\leanid{SameMPV₂Pos} directly; their proofs no longer convert an all-length
hypothesis back to positive length. The former bridge from positive-length
equality plus equal bond dimensions to an all-length statement has also been
removed from the public comparison surface. The empty-word coefficient lemma
remains only as an algebraic identity in the Lean source and is not presented
in the blueprint as a physical chain statement.

\section{Conclusion}

The block decomposition in
Eq.~\ref{eq:cpsv16_zero_tail_length_zero_decomposition_block_decomposition}
is correct as an intermediate linear-algebra statement. The former formal
defect concerned its shape:
Eq.~\ref{eq:cpsv16_zero_tail_length_zero_decomposition_all_length_identity}
combined a trivial empty-word dimension count with the substantive
positive-length comparison. The physical statement consists of
Eq.~\ref{eq:cpsv16_zero_tail_length_zero_decomposition_positive_length_comparison}
and the structural bound in
Eq.~\ref{eq:cpsv16_zero_tail_length_zero_decomposition_dimension_bound}.
The after-blocking primitive decomposition and all downstream consumers now
use this formulation.

\bibliographystyle{alpha}
\bibliography{references}

\end{document}
